\documentclass[12pt]{article}
\usepackage[utf8]{inputenc}
\usepackage[T1]{fontenc}
\usepackage{hyperref}
\title{Tunjung Language Resources for K-12 Curriculum Development}
\author{Compiled by Claude Sonnet 5 (claude-sonnet-5) for Andr\'as Kornai}
\date{}
\begin{document}
\maketitle
\section{Introduction}
Tunjung (ISO 639-3 tjg) is an Austronesian language of East Kalimantan, Indonesia. It is the only living member of the Mahakam branch. Sources give about 50,000 speakers, and the ethnic group is estimated at 62,000 to 76,000. Searches found no school materials in Tunjung, so little is documented for curriculum use.
\section{Primary Educational Resources}
No primers, textbooks or school programmes in Tunjung were found in the search results.
\section{Linguistic and Reference Materials}
\begin{itemize}
\item Wikipedia article on the Tunjung language, with classification, dialects and a wordlist drawn from Hudson (1967) via Blust: \url{https://en.wikipedia.org/wiki/Tunjung_language}
\item Numeral data for Tunjung, in the Max Planck EVA numerals database: \url{https://lingweb.eva.mpg.de/channumerals/Tunjung.htm}
\item Verbix language page for Tunjung: \url{https://docs.verbix.com/Languages/Tunjung}
\item A Universitas Mulawarman repository record (content not checked): \url{https://repository.unmul.ac.id/handle/123456789/37628}
\end{itemize}
\section{Religious and Community Resources}
\begin{itemize}
\item Wikipedia article on the Dayak Tunjung people: \url{https://en.wikipedia.org/wiki/Dayak_Tunjung_people}
\item Joshua Project profile of the Tunjung Dayak: \url{https://joshuaproject.net/people_groups/15643/ID}
\end{itemize}
\section{Digital and Audio Resources}
No audio, video or corpus resources were found in the search results.
\section{Curriculum Development Context}
Documentation is thin and no school use is known. A curriculum would first need fieldwork, an agreed orthography and local community partners. The dialects (Central Tunjung, Londong, Linggang, Berambai, Pahu) are listed differently by different sources.
\end{document}